\documentclass[11pt]{article}
\usepackage{fontspec}
\usepackage[french]{babel}
\usepackage{amsmath}
\usepackage{amssymb}
\usepackage{amsthm}
\usepackage{mathrsfs}
\usepackage{tikz-cd}
\usepackage{enumitem}

\newcommand{\printedpage}[1]{\marginpar{\footnotesize p.~#1}}
\newcommand{\msp}[1]{\marginpar{\footnotesize [#1]}}

\newcommand{\cE}{\mathcal{E}}
\newcommand{\cC}{\mathcal{C}}
\newcommand{\Ob}{\operatorname{Ob}}
\newcommand{\Fl}{\operatorname{Fl}}
\newcommand{\Hom}{\operatorname{Hom}}
\newcommand{\Ind}{\operatorname{Ind}}
\newcommand{\card}{\operatorname{card}}
\newcommand{\red}{\operatorname{\text{réd}}}
\newcommand{\id}{\mathrm{id}}
\newcommand{\Ens}{\mathrm{Ens}}
\newcommand{\circled}[1]{\tikz[baseline=(c.base)]\node[draw,circle,inner sep=1pt](c){#1};}

\begin{document}

\printedpage{164}
% head: A. Grothendieck : Dérivateurs, Ch. XVIII, Édition des Universités de Montpellier II et Paris VII

\msp{page 430}à prouver que pour tout $\lambda : s \longrightarrow t$, la flèche
\[
\varinjlim_{i} \Hom_{\lambda}(x, y_i) \ \longrightarrow \ \Hom_{\lambda}(x, y)
\]
est bijective. Pour $x$ et $\lambda$ fixés, et des systèmes inductifs variables dans $\cE_t$, grands devant $\pi$, cela signifie que le foncteur
\[
H_{\lambda}(x) \ = \ (y \longmapsto \Hom_{\lambda}(x, y)) \ : \ \cE_t \ \longrightarrow \ \Ens
\]
est $\pi$-accessible, et ceci quel que soit
\[
x \ \in \ \Ob \cE_{s\,\pi} \ .
\]
Or c'est le cas si $s \in \Ob \cE_{s\,\pi_0}$, en vertu de l'hypothèse $3^{\circ}$) (alors $H_{\lambda}(x_0)$ est $\pi_1$-accessible, % [sic] $s \in \Ob \cE_{s\,\pi_0}$ -> $x_0 \in \Ob \cE_{s\,\pi_0}$, the condition concerns an object of $\cE_{s\,\pi_0}$
et a fortiori $\pi$-accessible). Il faut passer de là au cas où $x \in \Ob \cE_{s\,\pi}$. Or comme $\pi_0$ est utile pour $\cE_s$ par $1^{\circ}$, on a
\[
x \ = \ \varinjlim_{j \in J} x_j \ , \qquad x_j \in \Ob \cE_{s\,\pi_0},\ J \text{ grande devant } \pi_0 \ ,
\]
et comme par $2^{\circ}$, $\pi^{\pi_0} = \pi \geqslant \pi_1 \geqslant \sup(\pi_0, \card \red \cE_{s\,\pi_0})$, on en conclut qu'on peut prendre $J$ de sorte que
\[
\card J \ \leqslant \ \pi \ .
\]
D'autre part, en vertu de $\pi \geqslant \pi_1 \ \geqslant \ \pi_{\lambda}$, le foncteur
\[
\left\{
\begin{array}{lcl}
x_0 & \longmapsto & H_{\lambda}(x_0) \\[1ex]
\cE_s & \longrightarrow & (\underline{\Hom}(\cE_s, \Ens))^{\circ}
\end{array}
\right.
\]
est $\pi_0$-accessible, donc on a
\[
H_{\lambda}(x_0) \ \simeq \ \varprojlim_{j \in J} H_{\lambda}(x_j) \ .
\]

\msp{page 431}Comme $x_j \ \in \ \Ob \cE_{s\,\pi_0}$, $H(x_j)$ est un foncteur $\pi_1$-accessible, et comme $\card J \leqslant \pi$, la $\varprojlim\limits_{j \in J} H_{\lambda}(x_j)$ est aussi $\pi$-accessible (cf. SGA 4 I 9.8).

Ainsi on a prouvé $\cE_{s\,\pi} \subset \cE(\pi)$ pour tout $s \in \Ob \cC$, i.e. l'inclusion
\[
\cE(\pi) \ \subset \ \cE_{\pi} \ ,
\]
d'où le foncteur pleinement fidèle
\[
\Ind_{\pi}(\cE(\pi)) \longrightarrow \cE \ , \tag*{$(*)$}
\]
comme $\cE$ satisfait à $\mathrm{L}_{\pi}$. Il reste à prouver que c'est essentiellement surjectif. Mais si $x \in \Ob \cE$, $s = f(x)$, on a $x \in \Ob \cE_s$, et comme $\pi_0$ est utile pour $\cE_s$ par $1^{\circ}$), et $\pi^{\pi_0} = \pi > \pi_0$,
$\pi$ est également utile pour $\pi$, donc on a \textit{dans} $\cE_s$ % [sic] utile pour $\pi$ -> utile pour $\cE_s$, usefulness is a property relative to the category $\cE_s$
\[
x \ = \ \varinjlim_{I} x_i \ , \qquad \text{les } x_i \in \cE_{s\,\pi} \subset \cE(\pi),\ I \text{ grande devant } \pi \ ,
\]
\printedpage{165}
% head: A. Grothendieck : Dérivateurs, Ch. XVIII, Édition des Universités de Montpellier II et Paris VII
et d'après a) (comme $\pi \geqslant \pi_1 \geqslant \sup_{\lambda \in \Fl \cC} \pi_{\lambda}$) $x$ est aussi limite des $x_i$ dans $\cE$, donc il est dans l'image essentielle de $(*)$.

d) Reste à prouver que l'on a $\cE_{\pi} \ \subset \ \cE(\pi)$. Mais on sait (par 4.8.~\dots), comme $\cE \ \simeq \ \Ind_{\pi}(\cE(\pi))$, que $\cE_{\pi} \ = \ \cE(\pi)^{\natural}$, i.e. que pour $x \ \in \ \Ob \cE_{\pi}$, $\exists \ y \ \in \ \cE(\pi)$ et un projecteur $p$ dans $x$, tel que $x$ soit isomorphe à l'image de $p$ dans $\cE$. Or $f(p)$ est un projecteur dans $f(y) = s$, donc par hypothèse sur $\cC$, $f(p) = \id_s$, i.e. $p$ est un projecteur dans $\cE_s$. Comme donc $x$ est dans $\cE_s$, et

\msp{page 432}même dans $\cE_{s\,\pi}$, puisque $\cE_{s\,\pi}$ est stable dans $\cE_s$ par facteurs directs. Cela prouve que $x \in \Ob \cE(\pi)$, q.e.d.

On peut présenter 4.15.1 sous une forme plus sympathique.

\textbf{Corollaire 4.15.4.} \textit{Soit $f : \cE \longrightarrow \cC$ un foncteur de $\mathfrak{U}$-catégories. Les conditions suivantes sont équivalentes :}
\begin{enumerate}
\item[\textit{a)}] \textit{Pour tout foncteur $\cC' \longrightarrow \cC$, avec $\cC'$ petite et sans projecteurs, $\cE' \overset{\text{déf}}{=} \cE \times_{\cC} \cC'$ est accessible, et $f' : \cE' \longrightarrow \cC'$ est accessible.}
\item[\textit{b)}] \textit{Comme a), mais avec $\cC' = \mathbf{\Delta}_1$, en d'autres termes, c'est la condition $H$ p.~409 : pour toute $\lambda \in \Fl \cC$, $\cE_{\lambda} = \cE \times_{\cC} (\mathbf{\Delta}_1, i_{\lambda})$ est accessible, et $\cE_{\lambda} \longrightarrow \mathbf{\Delta}_1$ est accessible.}
\item[\textit{c)}] \textit{La condition $H' = H\ \mathit{1} + H\ \mathit{2} + H\ \mathit{3}$, explicitée dans l'énoncé de 4.15.1.}
\end{enumerate}
\textsc{Démonstration}. Elle se fait circulaire,
\[
\mathrm{a} \ \Longrightarrow \ \mathrm{b} \ \Longrightarrow \ \mathrm{c} \ \Longrightarrow \ \mathrm{a} \ ,
\]
l'implication a $\Longrightarrow$ b étant tautologique, b $\Longrightarrow$ c l'étant presque (cf. 4.14.24, a $\Longrightarrow$ b $\Longrightarrow$ c). En fait, les conditions b et c sont des conditions sur les catégories $\cE_s \xrightarrow{f_{\lambda}} \mathbf{\Delta}_1$ pour $\lambda \in \Fl \cC$, donc stables par changement de base. Pour prouver

\msp{page 433}c $\Longrightarrow$ a, on peut donc supposer que $\cC' = \cC$, $\cC' \longrightarrow \cC$ le foncteur identique, et prouver que si c) est satisfaite et si $\cC$ est petite et sans projecteurs, alors $\cE$ est accessible et $f$ accessible. Or c'est ce qu'affirme 4.15.1.

\textbf{Remarque 4.15.5.} Si $f : \cE \longrightarrow \cC$ est un foncteur pleinement fidèle, il est immédiat que la condition H est satisfaite (car $\cE_{\lambda} \longrightarrow \mathbf{\Delta}_1$ sera pleinement fidèle, donc $\cE_{\lambda}$ est équivalente
à $e$ où à $\mathbf{\Delta}_1$, donc accessible, et tout foncteur de $\cE_{\lambda}$ dans une catégorie quelconque est % [sic] où -> ou, disjunction "e ou Δ1"
accessible \dots). Mais si on suppose $\cC$ petite \textit{sans} supposer $\cC$ sans projecteurs, il ne s'ensuit pas que $\cE$ soit accessible, ou ce qui revient au même (étant essentiellement petite), que $\cE$ soit karoubienne. (L'exemple le plus trivial est celui où $\cE = \cC$, $f = \id_{\cC}$, $\cE$ n'est karoubienne que si $\cC$ l'est !) Mais même en supposant $\cC$ une petite catégorie karoubienne, il ne s'ensuit pas forcément que $\cE$ l'est. En effet, \textit{toute} petite catégorie $\cE$ (karoubienne ou non) se plonge comme sous-catégorie pleine dans son enveloppe karoubienne $\cC$ -- il suffit donc de prendre $\cE$ non karoubienne. L'exemple

\msp{page 434}le plus économique est
\[
\Ob \cE \ = \ \{0\} \ , \qquad \Fl(\cE) \ = \ \{\id_0, p\} \ , \qquad \text{avec } p^2 = p \neq \id_0 \ ,
\]
\printedpage{166}
% head: A. Grothendieck : Dérivateurs, Ch. XVIII, Édition des Universités de Montpellier II et Paris VII
l'enveloppe de Karoubi de $\cE$ est donnée par
\[
P \ = \ \left\{
\begin{tikzcd}
1 \arrow[r, shift right, "\sigma"'] & 0 \arrow[l, shift right, "\rho"'] \arrow[loop right, "p"]
\end{tikzcd}
\qquad
\begin{aligned}
\rho\sigma \ &= \ \id_1 \\
\sigma\rho \ &= \ p \ \neq \ \id_0
\end{aligned}
\quad \Longrightarrow \ p^2 = p
\right\} .
\]
Cela soulève la question de trouver un énoncé similaire à 4.15.1, donnant des conditions serviables simples, plus fortes que $\mathrm{H}' = \mathrm{H}$ (1+2+3), sur $f : \cE \longrightarrow \cC$, pour que $\cE$ soit accessible et $f$ accessible, en supposant $\cC$ \textit{karoubienne}. Ou encore, des conditions nécessaires et
suffisants serviables (plus fortes que 4.15.4) sur $f : \cE \longrightarrow \cC$, pour que toute \textit{petite} catégorie % [sic] suffisants -> suffisantes, agreement with "conditions"
\textit{karoubienne} $\cC'$, et tout foncteur $\cC' \longrightarrow \cC$, la catégorie $\cE' = \cE \times_{\cC} \cC'$ soit accessible, et le foncteur $\cE' \longrightarrow \cC'$ accessible. Une idée naturelle, c'est qu'il suffirait que cette conclusion soit vraie quand $\cC'$ est une des deux catégories $\mathbf{\Delta}_1$, $P$ (où $P$ est la catégorie à deux objets, à trois flèches non identiques, définie plus haut). Mais je doute qu'il en soit ainsi (\footnote[178]{Ça devrait marcher par contre si on prend les catégories $0 \longrightarrow 1$,
\[
\begin{array}{c}
2 \longrightarrow \underbrace{\begin{tikzcd}[ampersand replacement=\&]
1 \arrow[r, shift right, "\sigma"'] \& 0 \arrow[l, shift right, "\rho"'] \arrow[loop right, "p"]
\end{tikzcd}}_{P}
\quad \text{et} \quad
2 \longleftarrow \underbrace{\begin{tikzcd}[ampersand replacement=\&]
1 \arrow[r, shift right, "\sigma"'] \& 0 \arrow[l, shift right, "\rho"'] \arrow[loop right, "p"]
\end{tikzcd}}_{P}
\qquad (\text{\footnotemark[179]}) \ , \\
\nwarrow \ \begin{array}{c} \text{3 objets et} \\ \text{5 flèches non} \\ \text{identiques} \end{array} \ \nearrow
\end{array}
\]
grâce à 4.15.7, page 440.}\footnotetext[179]{ces catégories représentent respectivement les foncteurs $\mathrm{Cat} \longrightarrow \Ens$ associant à une catégorie $\cC$ les diagrammes dans $\cC$ formés de $(x, y, p, u)$, avec $p$ projecteur dans $p$ \texttt{[plutôt dans $\cC$]}, et $u : y \longrightarrow x^{\natural}$ (i.e. $u : y \longrightarrow x$ avec $u = pu$), resp. $u : x^{\natural} \longrightarrow y$ (i.e. $u : x \longrightarrow y$ avec $u = up$).}).

On va quand même essayer un énoncé un peu très particulier dans ce sens :

% ?? uncertain: hand-drawn sign in the left margin beside items a) and b) of Proposition 4.15.6 -- a curved mark resembling a bracket or the figure "2"; meaning unknown, not transcribed as text
\msp{page 435}\textbf{Proposition 4.15.6} (\footnote[180]{\textit{Faux ?}}). \textit{Soit $f : \cE \longrightarrow \cC$ un foncteur de $\mathfrak{U}$-catégories. On suppose :}
\begin{enumerate}
\item[\textit{a)}] \textit{$f$ est fibrant ou cofibrant} (\footnote[181]{Le cas fibrant n'est pas prouvé, peut-être \textit{faux} -- et de même pour le cas cofibrant !}).
\item[\textit{b)}] \textit{Les catégories fibres $\cE_s$ ($s \in \Ob \cC$) sont accessibles, et pour toute flèche $\lambda : s \longrightarrow t$ dans $\cC$, le foncteur changement de base $\lambda^* : \cE_t \longrightarrow \cE_s$ (resp. cochangement de base $\lambda_* : \cE_s \longrightarrow \cE_t$) est accessible.}
\item[\textit{c)}] \textit{$\cC$ est essentiellement petite et karoubienne.}
\end{enumerate}
\textit{Alors $\cE$ est accessible, et $f$ est accessible.}

\textsc{Démonstration}. On a déjà vu que les conditions a) et b) impliquent $\mathrm{H}'$, i.e. H (1+2+3). D'ailleurs, si $\cC'$ est une catégorie et $\cC' \xrightarrow{\ \sim\ } \cC$ une équivalence, donc alors $\cE' \overset{\text{déf}}{=} \cE \times_{\cC} \cC' \longrightarrow \cE$ est une équivalence, donc la conclusion voulue pour $\cE$, $f$ équivaut à la même conclusion pour $\cE'$, $f' : \cE' \longrightarrow \cC'$. Donc on peut supposer que $\cC$ est petite (et même réduite, si on y tient). En vertu de 4.15.2 a), on voit donc que les foncteurs d'inclusion $\cE_s \longrightarrow \cE$ sont accessibles (et même $\pi_1$-accessibles) dès que $\pi_1 \geqslant \sup_{\lambda \in \Fl \cC} \pi_{\lambda}$. On voit d'ailleurs

\msp{page 436}par AQT \texttt{[âne qui trotte]} qu'on peut trouver des cardinaux $\pi_0$, $\pi_1$ ayant les propriétés suivantes :

\printedpage{167}
% head: A. Grothendieck : Dérivateurs, Ch. XVIII, Édition des Universités de Montpellier II et Paris VII
\begin{enumerate}
\item[\circled{1\textsuperscript{o}}] $\pi_0$ est utile pour $\cE_s$, $\forall \ s \in \Ob \cC$.
\item[$2^{\circ}$)] $\pi_1 \geqslant \sup(\pi_0, \card \cC, \sup_{\lambda \in \Fl \cC} \pi_{\lambda}, \sup_{s \in \Ob \cC} \card \red \cE_s)$.
\item[\circled{3\textsuperscript{o}}] $\forall \ (\lambda : s \longrightarrow t) \in \Fl \cC$ et $x_0 \in \Ob \cE_{s\,\pi_0}$, le foncteur
\[
H_{\lambda}(x_0) \ : \ (x_1 \longmapsto \Hom_{\lambda}(x_0, x_1)) \ : \ \cE_t \ \longrightarrow \ (\Ens)
\]
est $\pi_1$-accessible. (\textbf{N.B.} Si $f$ est cofibrant, cela signifie que $\forall \ \lambda \in \Fl \cC$, $\lambda : s \longrightarrow t$,
\fbox{$\lambda_*(\cE_{s\,\pi_0}) \ \subset \ \cE_{s\,\pi_1}$}. Si $f$ est fibrant, il suffit d'exiger que $\forall \ \lambda \in \Fl \cC$, $\lambda^* : \cE_t \longrightarrow \cE_s$ % [sic] $\cE_{s\,\pi_1}$ -> $\cE_{t\,\pi_1}$, $\lambda_*$ goes to $\cE_t$, cf. the later display "(contenu dans 3°)"
soit $\pi_1$-accessible.) Mais dans ce cas, on veut aussi \fbox{$\lambda^*(\cE_{s\,\pi_0}) \ \subset \ \cE_{s\,\pi_1}$}. % [sic] $\lambda^*(\cE_{s\,\pi_0})$ -> $\lambda^*(\cE_{t\,\pi_0})$, $\lambda^*$ is defined on $\cE_t$, cf. the later display "(contenu dans 3°)"
\item[\circled{4\textsuperscript{o}}] Les foncteurs $\lambda^*$ (cas $f$ fibrant) resp. $\lambda_*$ (cas $f$ cofibrant) sont $\pi_1$-accessibles.
\end{enumerate}
Soit alors $\pi$ un cardinal tel que
\[
\pi \ = \ \pi^{\pi_0} \ \geqslant \ \pi_1 \ , \tag*{$(*)$}
\]
et soit $\cE(\pi)$ la sous-catégorie (strictement) pleine de $\cE$ définie par
\[
\Ob \cE(\pi) \ = \ \bigcup_{s \, \in \, \Ob \cC} \Ob \cE_{s\,\pi} \ .
\]
Grâce à $(*)$, et à $1^{\circ}$, $4^{\circ}$ et à
\[
\left\{
\begin{array}{l}
\pi_1 \ \geqslant \ \sup_s(\card \red \cE_{s\,\pi_0}) \qquad \text{(contenu dans $2^{\circ}$)} \\[1ex]
\lambda^*(\cE_{t\,\pi_0}) \ \subset \ \cE_{s\,\pi_1} \qquad \text{resp.} \qquad \lambda_*(\cE_{s\,\pi_0}) \ \subset \ \cE_{t\,\pi_1} \qquad \text{(contenu dans $3^{\circ}$)} \ ,
\end{array}
\right.
\]
on voit que $\cE(\pi)$ est une sous-catégorie

\msp{page 437}fibrée, resp. une sous-catégorie cofibrée, sur $\cC$ (le foncteur $\cE(\pi) \longrightarrow \cE$ étant donc cartésien resp. cocartésien). À $\cC$-équivalence près, $\cE$ se reconstitue alors à partir de $\cE(\pi)$, en passant au $\Ind_{\pi}(\cE_{s\,\pi})$, qui forment un système fibré resp. cofibré sur $\cC$ par les foncteurs $\Ind_{\pi}(\lambda^*_{\pi})$ resp. $\Ind_{\pi}(\lambda^{\pi}_*)$ \dots

Prouvons que $\cE$ satisfait à $\mathrm{L}_{\pi}$ et que $f : \cE \longrightarrow \cC$ est $\pi$-accessible (\footnote[182]{Ici on utilise seulement $\pi \geqslant \card \red \cC$.}). Comme $\cC$ est karoubienne, et $\pi \geqslant \pi_1 \geqslant \card \cC$, $\cC$ satisfait à $\mathrm{L}_{\pi}$ et les ind-objets grands devant $\pi$ de $\cC$ sont essentiellement constants (4.8.13). Soit alors $I$ grande devant $\pi$, et $I \longrightarrow \cE$, $i \longmapsto x_i$, un foncteur, il faut construire une $\varinjlim$ dans $\cE$, en utilisant la $\varinjlim$ déjà connue $s$ des $s_i = f(x_i)$. On va prendre un $x$ dans $\cE_s$. Prenons d'abord le cas où $f$ est cofibrant. On voit que les $\beta_i(x_i)$ forment un système inductif dans $\cE_s$, puisque pour $i = j$, on a
\[
\beta_{i\,*}(x_j) \ = \ \beta_{j\,*} \underbrace{\beta_{ji\,*}(x_i)}_{\longrightarrow \, x_j} \ \longrightarrow \ \beta_{j\,*}(x_j) \ . % [sic] $\beta_{i\,*}(x_j)$ -> $\beta_{i\,*}(x_i)$, the object transported from $s_i$ is $x_i$
\]
\[
\begin{tikzcd}[row sep=small]
x_i \arrow[rr, "\alpha_i"] & & x \, ? \\
s_i \arrow[rr, "\beta_i"] \arrow[dr, "\beta_{ji}"'] & & s \\
& s_j \arrow[ur, "\beta_j"'] &
\end{tikzcd}
\]
On pose

\printedpage{168}
% head: A. Grothendieck : Dérivateurs, Ch. XVIII, Édition des Universités de Montpellier II et Paris VII
\msp{page 438}
\[
x \ = \ \varinjlim_{I} \beta_{i\,*}(x_i) \ , \qquad \varinjlim \ \textit{dans } \cE_s \ .
\]
On sait que c'est une $\varinjlim$ dans $\cE$, donc que pour tout $y$ dans $\cE$, on a
\[
\Hom(x, y) \xrightarrow{\ \sim\ } \varprojlim_{I} \Hom(\beta_{i\,*}(x_i), y) \ ,
\]
ou ce qui revient au même, pour toute flèche $\lambda : s \longrightarrow t$, et $y$ dans $\Ob \cE_t$, on a
\[
\underbrace{\Hom_{\lambda}(x, y)}_{= \, \Hom_{\cE_t}(\lambda_* x, y)} \ \simeq \ \varprojlim \Hom_{\cE_t}(\underbrace{\lambda_* \beta_{i\,*}(x_i)}_{(\lambda_i \beta_i)_*(x_i)}, y) \ .
\]
Cela résulte aussi du fait que $\lambda_*$ est $\pi$-accessible (étant $\pi_1$-accessible) \dots{} On vérifie facilement que c'est une $\varinjlim$ des $x_i$ dans $\cE$. Et bien sûr, cette limite commute avec $f$ par construction.

Mais dans le cas où $f$ est fibrant, je ne vois comment faire pour construire une limite dans $\cE_s$. Il est vrai qu'on peut supposer que tous les $s_i$ sont égaux à $s$, mais les transitions $\beta_{ji}$ ne sont pas l'identité, ni même des isomorphismes. Il faudrait prendre $\cC = P$ pour essayer de comprendre -- je ne vais pas insister (\footnote[183]{J'ai vérifié que si $f$ est fibrant, alors $\cE$ est \textit{karoubienne}. Il y a donc de l'espoir \dots}).

\msp{page 439}Dans le cas où $f$ est fibrant, admettons (comme hypothèse supplémentaire au besoin) que $\cE$ satisfasse $L$ et que $f$ soit accessible, et prenons $\pi$ assez grand pour que $\cE$, $f$ satisfassent $\mathrm{L}_{\pi}$. Nous pouvons alors terminer la démonstration de 4.15.6. Il faut prouver
\[
\cE(\pi) \ \subset \ \cE_{\pi} \ ,
\]
donc que si $x \in \Ob \cE_{s\,\pi}$, et $(y_i)_{i \in I}$ est un système inductif dans $\cE$ grand devant $\pi$, de limite $y$, on a
\[
\varinjlim \Hom_{\cE}(x, y_i) \xrightarrow{\ \sim\ } \Hom_{\cE}(x, y) \ .
\]
Donc il faut prouver que toute flèche $x \xrightarrow{u} y$ se factorise, de façon essentiellement unique $u_i$, par un $y_i$. Mais la flèche $s \xrightarrow{\ \lambda \, = \, f(u)\ } t$ dans $\cC$ se factorise, de façon essentiellement unique, par des $\lambda_i \, : \, s \longrightarrow t_i$. Donc on est ramené à trouver un \textit{$\lambda_i$-morphisme} $x \longrightarrow y_i$ pour $i$ grand, qui factorise. Dans le cas fibrant, cela signifie qu'on cherche $x \longrightarrow \lambda^*_i(y_i)$ qui factorise $x \longrightarrow \lambda^*(y)$, or les $\lambda^*_i(y_i)$ forment un système inductif dans $\cE_s$, et sa limite \textit{dans} $\cE_s$ est $\lambda^*(y)$ (\textit{à vérifier}). Comme $x$ est dans $\cE_{s\,\pi}$, l'hom\texttt{[omorphisme]} donné $x \longrightarrow \lambda^*(y)$ se factorise de

\msp{page 440}façon essentiellement unique par un $\lambda^*_i(y_i)$.

Cette fois par contre, c'est le cas $f$ cofibrant qui pose problème. En fait, je n'arrive pas à m'en tirer.

Ainsi, on trouve bien un foncteur
\[
\Ind_{\pi}(\cE(\pi)) \ \longrightarrow \ \cE \ ,
\]

\end{document}
